% Options for packages loaded elsewhere
\PassOptionsToPackage{unicode}{hyperref}
\PassOptionsToPackage{hyphens}{url}
\documentclass[
]{article}
\usepackage{xcolor}
\usepackage[margin=1in]{geometry}
\usepackage{amsmath,amssymb}
\setcounter{secnumdepth}{5}
\usepackage{iftex}
\ifPDFTeX
  \usepackage[T1]{fontenc}
  \usepackage[utf8]{inputenc}
  \usepackage{textcomp} % provide euro and other symbols
\else % if luatex or xetex
  \usepackage{unicode-math} % this also loads fontspec
  \defaultfontfeatures{Scale=MatchLowercase}
  \defaultfontfeatures[\rmfamily]{Ligatures=TeX,Scale=1}
\fi
\usepackage{lmodern}
\ifPDFTeX\else
  % xetex/luatex font selection
\fi
% Use upquote if available, for straight quotes in verbatim environments
\IfFileExists{upquote.sty}{\usepackage{upquote}}{}
\IfFileExists{microtype.sty}{% use microtype if available
  \usepackage[]{microtype}
  \UseMicrotypeSet[protrusion]{basicmath} % disable protrusion for tt fonts
}{}
\makeatletter
\@ifundefined{KOMAClassName}{% if non-KOMA class
  \IfFileExists{parskip.sty}{%
    \usepackage{parskip}
  }{% else
    \setlength{\parindent}{0pt}
    \setlength{\parskip}{6pt plus 2pt minus 1pt}}
}{% if KOMA class
  \KOMAoptions{parskip=half}}
\makeatother
\usepackage{color}
\usepackage{fancyvrb}
\newcommand{\VerbBar}{|}
\newcommand{\VERB}{\Verb[commandchars=\\\{\}]}
\DefineVerbatimEnvironment{Highlighting}{Verbatim}{commandchars=\\\{\}}
% Add ',fontsize=\small' for more characters per line
\usepackage{framed}
\definecolor{shadecolor}{RGB}{248,248,248}
\newenvironment{Shaded}{\begin{snugshade}}{\end{snugshade}}
\newcommand{\AlertTok}[1]{\textcolor[rgb]{0.94,0.16,0.16}{#1}}
\newcommand{\AnnotationTok}[1]{\textcolor[rgb]{0.56,0.35,0.01}{\textbf{\textit{#1}}}}
\newcommand{\AttributeTok}[1]{\textcolor[rgb]{0.13,0.29,0.53}{#1}}
\newcommand{\BaseNTok}[1]{\textcolor[rgb]{0.00,0.00,0.81}{#1}}
\newcommand{\BuiltInTok}[1]{#1}
\newcommand{\CharTok}[1]{\textcolor[rgb]{0.31,0.60,0.02}{#1}}
\newcommand{\CommentTok}[1]{\textcolor[rgb]{0.56,0.35,0.01}{\textit{#1}}}
\newcommand{\CommentVarTok}[1]{\textcolor[rgb]{0.56,0.35,0.01}{\textbf{\textit{#1}}}}
\newcommand{\ConstantTok}[1]{\textcolor[rgb]{0.56,0.35,0.01}{#1}}
\newcommand{\ControlFlowTok}[1]{\textcolor[rgb]{0.13,0.29,0.53}{\textbf{#1}}}
\newcommand{\DataTypeTok}[1]{\textcolor[rgb]{0.13,0.29,0.53}{#1}}
\newcommand{\DecValTok}[1]{\textcolor[rgb]{0.00,0.00,0.81}{#1}}
\newcommand{\DocumentationTok}[1]{\textcolor[rgb]{0.56,0.35,0.01}{\textbf{\textit{#1}}}}
\newcommand{\ErrorTok}[1]{\textcolor[rgb]{0.64,0.00,0.00}{\textbf{#1}}}
\newcommand{\ExtensionTok}[1]{#1}
\newcommand{\FloatTok}[1]{\textcolor[rgb]{0.00,0.00,0.81}{#1}}
\newcommand{\FunctionTok}[1]{\textcolor[rgb]{0.13,0.29,0.53}{\textbf{#1}}}
\newcommand{\ImportTok}[1]{#1}
\newcommand{\InformationTok}[1]{\textcolor[rgb]{0.56,0.35,0.01}{\textbf{\textit{#1}}}}
\newcommand{\KeywordTok}[1]{\textcolor[rgb]{0.13,0.29,0.53}{\textbf{#1}}}
\newcommand{\NormalTok}[1]{#1}
\newcommand{\OperatorTok}[1]{\textcolor[rgb]{0.81,0.36,0.00}{\textbf{#1}}}
\newcommand{\OtherTok}[1]{\textcolor[rgb]{0.56,0.35,0.01}{#1}}
\newcommand{\PreprocessorTok}[1]{\textcolor[rgb]{0.56,0.35,0.01}{\textit{#1}}}
\newcommand{\RegionMarkerTok}[1]{#1}
\newcommand{\SpecialCharTok}[1]{\textcolor[rgb]{0.81,0.36,0.00}{\textbf{#1}}}
\newcommand{\SpecialStringTok}[1]{\textcolor[rgb]{0.31,0.60,0.02}{#1}}
\newcommand{\StringTok}[1]{\textcolor[rgb]{0.31,0.60,0.02}{#1}}
\newcommand{\VariableTok}[1]{\textcolor[rgb]{0.00,0.00,0.00}{#1}}
\newcommand{\VerbatimStringTok}[1]{\textcolor[rgb]{0.31,0.60,0.02}{#1}}
\newcommand{\WarningTok}[1]{\textcolor[rgb]{0.56,0.35,0.01}{\textbf{\textit{#1}}}}
\usepackage{longtable,booktabs,array}
\newcounter{none} % for unnumbered tables
\usepackage{calc} % for calculating minipage widths
% Correct order of tables after \paragraph or \subparagraph
\usepackage{etoolbox}
\makeatletter
\patchcmd\longtable{\par}{\if@noskipsec\mbox{}\fi\par}{}{}
\makeatother
% Allow footnotes in longtable head/foot
\IfFileExists{footnotehyper.sty}{\usepackage{footnotehyper}}{\usepackage{footnote}}
\makesavenoteenv{longtable}
\usepackage{graphicx}
\makeatletter
\newsavebox\pandoc@box
\newcommand*\pandocbounded[1]{% scales image to fit in text height/width
  \sbox\pandoc@box{#1}%
  \Gscale@div\@tempa{\textheight}{\dimexpr\ht\pandoc@box+\dp\pandoc@box\relax}%
  \Gscale@div\@tempb{\linewidth}{\wd\pandoc@box}%
  \ifdim\@tempb\p@<\@tempa\p@\let\@tempa\@tempb\fi% select the smaller of both
  \ifdim\@tempa\p@<\p@\scalebox{\@tempa}{\usebox\pandoc@box}%
  \else\usebox{\pandoc@box}%
  \fi%
}
% Set default figure placement to htbp
\def\fps@figure{htbp}
\makeatother
\setlength{\emergencystretch}{3em} % prevent overfull lines
\providecommand{\tightlist}{%
  \setlength{\itemsep}{0pt}\setlength{\parskip}{0pt}}
\usepackage{bookmark}
\IfFileExists{xurl.sty}{\usepackage{xurl}}{} % add URL line breaks if available
\urlstyle{same}
\hypersetup{
  pdftitle={Quantifying a Social Problem: Housing Inequality in the Netherlands},
  pdfauthor={Group CPR02: Joey Pekel, Lior Santcroos, Benjamin Lobatto, Jeroen Warnaar, Steven van Zuijlen, Bianca Pariamachi},
  hidelinks,
  pdfcreator={LaTeX via pandoc}}

\title{Quantifying a Social Problem: Housing Inequality in the
Netherlands}
\usepackage{etoolbox}
\makeatletter
\providecommand{\subtitle}[1]{% add subtitle to \maketitle
  \apptocmd{\@title}{\par {\large #1 \par}}{}{}
}
\makeatother
\subtitle{Affordability, regional disparity, and the price-to-income
ratio}
\author{Group CPR02: Joey Pekel, Lior Santcroos, Benjamin Lobatto,
Jeroen Warnaar, Steven van Zuijlen, Bianca Pariamachi}
\date{17 June 2026}

\begin{document}
\maketitle

{
\setcounter{tocdepth}{2}
\tableofcontents
}
\begin{quote}
\textbf{Public GitHub repository:}
\url{https://github.com/benjaminlobatto/woningmarkt-project} \emph{(This
link is required and must be visible in the handed-in PDF.)}
\end{quote}

\section{Introduction and context}\label{introduction-and-context}

Housing affordability is one of the most pressing social problems in the
Netherlands. Over the past two decades, house prices have risen far
faster than household incomes, and the gap is highly uneven across
provinces: the Randstad (Noord-Holland, Zuid-Holland, Utrecht) has
decoupled from peripheral regions. The OECD and UN-Habitat both treat
the price-to-income ratio (PIR) as a core affordability indicator, with
a value above 3 widely read as ``unaffordable'' (UN-Habitat, n.d.). The
OECD's 2025 Economic Survey of the Netherlands documents a structural
shortage of roughly 400,000 homes and a three-tiered market that
increasingly squeezes lower- and middle-income households (OECD, 2025).
De Nederlandsche Bank reaches a similar conclusion from the demand side:
house prices have outpaced households' borrowing capacity since 2013, so
that by 2027 only about one in three households is projected to be able
to finance an average home, down from nearly one in two in 2019 (DNB,
2024). These two independent assessments --- one structural, one
financial --- motivate our quantitative analysis below.

This decoupling matters because it shifts wealth toward existing owners
and locks younger and lower-income households out of ownership,
reinforcing inter-generational inequality. We quantify the problem along
three rubric dimensions --- over time, across provinces (space), and
around a specific event (the 2008 financial crisis) --- and across an
income sub-population split.

\subsection{Research questions}\label{research-questions}

\begin{enumerate}
\def\labelenumi{\arabic{enumi}.}
\tightlist
\item
  \textbf{Temporal:} How has national housing affordability (PIR)
  evolved over time?
\item
  \textbf{Spatial:} Which provinces are least affordable, and how large
  is the gap?
\item
  \textbf{Event:} How did the 2008 crisis interrupt the price
  trajectory?
\item
  \textbf{Sub-population:} Do affordability dynamics differ between
  higher- and lower-income provinces?
\end{enumerate}

\section{Data and methods}\label{data-and-methods}

We document and merge two public CBS datasets, both retrieved
programmatically via the \texttt{cbsodataR} OData API (collector:
Statistics Netherlands; method: administrative registers and the Land
Registry \emph{Kadaster}).

{\def\LTcaptype{none} % do not increment counter
\begin{longtable}[]{@{}
  >{\raggedright\arraybackslash}p{(\linewidth - 6\tabcolsep) * \real{0.2500}}
  >{\raggedright\arraybackslash}p{(\linewidth - 6\tabcolsep) * \real{0.2500}}
  >{\raggedright\arraybackslash}p{(\linewidth - 6\tabcolsep) * \real{0.2500}}
  >{\raggedright\arraybackslash}p{(\linewidth - 6\tabcolsep) * \real{0.2500}}@{}}
\toprule\noalign{}
\begin{minipage}[b]{\linewidth}\raggedright
Dataset
\end{minipage} & \begin{minipage}[b]{\linewidth}\raggedright
What it provides
\end{minipage} & \begin{minipage}[b]{\linewidth}\raggedright
Granularity
\end{minipage} & \begin{minipage}[b]{\linewidth}\raggedright
CBS table
\end{minipage} \\
\midrule\noalign{}
\endhead
\bottomrule\noalign{}
\endlastfoot
House prices (existing own homes) & Average sale price, € & Year ×
province & \texttt{83625NED} \\
Household income & Median disposable income, € & Year × province &
\texttt{86161NED} \\
\end{longtable}
}

We chose CBS because it is the authoritative national statistical
office, its tables are openly licensed and versioned, and both
indicators share a common \texttt{year\ ×\ province} key that lets us
merge them and compute affordability directly. The price series runs
from 1995, while the regional income series begins in 2011; the merged
affordability panel therefore covers 2011--2024, and the 2008 event
analysis uses the longer price series. A limitation is that purchase
prices ignore the large rental/social sector (discussed in §6).

\begin{Shaded}
\begin{Highlighting}[]
\CommentTok{\# Verified province{-}level CBS tables (confirmed via cbs\_get\_meta to have a}
\CommentTok{\# RegioS province dimension and yearly Perioden). Each download is cached to}
\CommentTok{\# data/ on first knit.}
\NormalTok{raw\_price  }\OtherTok{\textless{}{-}} \FunctionTok{try\_cbs}\NormalTok{(}\StringTok{"83625NED"}\NormalTok{)   }\CommentTok{\# avg sale price, existing homes, by province}
\NormalTok{raw\_income }\OtherTok{\textless{}{-}} \FunctionTok{try\_cbs}\NormalTok{(}\StringTok{"86161NED"}\NormalTok{)   }\CommentTok{\# disposable household income, by province}
\end{Highlighting}
\end{Shaded}

The data source for this knit is: r data\_source`.

\section{Data cleaning and merging}\label{data-cleaning-and-merging}

\texttt{normalise\_cbs()} trims CBS' padded codes, parses the period
into a numeric \texttt{year}, restricts to the 12 provinces (codes
\texttt{PV20}--\texttt{PV31}), and rescales income from CBS' ``× 1000
euro'' unit. We then merge the two datasets on
\texttt{province\ ×\ year} and construct two new analytical variables:

\begin{itemize}
\tightlist
\item
  \textbf{\texttt{pir}} --- the price-to-income ratio
  (\texttt{avg\_price\ /\ median\_income}), the affordability indicator;
  and
\item
  \textbf{\texttt{gap}} --- the affordability gap index: PIR relative to
  the UN-Habitat threshold of 3 (\texttt{pir/3\ -\ 1}), so positive
  values flag unaffordable province-years and the magnitude is directly
  interpretable.
\end{itemize}

\begin{Shaded}
\begin{Highlighting}[]
\NormalTok{panel }\OtherTok{\textless{}{-}} \FunctionTok{inner\_join}\NormalTok{(prices, incomes, }\AttributeTok{by =} \FunctionTok{c}\NormalTok{(}\StringTok{"province"}\NormalTok{,}\StringTok{"year"}\NormalTok{)) }\SpecialCharTok{|\textgreater{}}
  \FunctionTok{filter}\NormalTok{(province }\SpecialCharTok{\%in\%}\NormalTok{ provinces) }\SpecialCharTok{|\textgreater{}}
  \FunctionTok{mutate}\NormalTok{(}
    \AttributeTok{pir =} \FunctionTok{price\_to\_income}\NormalTok{(avg\_price, median\_income),  }\CommentTok{\# new variable 1}
    \AttributeTok{gap =}\NormalTok{ pir}\SpecialCharTok{/}\DecValTok{3} \SpecialCharTok{{-}} \DecValTok{1}                                    \CommentTok{\# new variable 2}
\NormalTok{  ) }\SpecialCharTok{|\textgreater{}}
  \FunctionTok{arrange}\NormalTok{(province, year)}

\CommentTok{\# National yearly aggregates (unweighted mean across provinces).}
\CommentTok{\# PIR uses the merged panel (2011{-}2024); the price{-}only series spans 1995{-}2024}
\CommentTok{\# and is used for the 2008 event analysis.}
\NormalTok{national }\OtherTok{\textless{}{-}}\NormalTok{ panel }\SpecialCharTok{|\textgreater{}}
  \FunctionTok{group\_by}\NormalTok{(year) }\SpecialCharTok{|\textgreater{}}
  \FunctionTok{summarise}\NormalTok{(}\AttributeTok{avg\_price =} \FunctionTok{mean}\NormalTok{(avg\_price), }\AttributeTok{median\_income =} \FunctionTok{mean}\NormalTok{(median\_income),}
            \AttributeTok{pir =} \FunctionTok{mean}\NormalTok{(pir), }\AttributeTok{.groups =} \StringTok{"drop"}\NormalTok{)}
\NormalTok{price\_history }\OtherTok{\textless{}{-}}\NormalTok{ prices }\SpecialCharTok{|\textgreater{}}
  \FunctionTok{group\_by}\NormalTok{(year) }\SpecialCharTok{|\textgreater{}} \FunctionTok{summarise}\NormalTok{(}\AttributeTok{avg\_price =} \FunctionTok{mean}\NormalTok{(avg\_price), }\AttributeTok{.groups =} \StringTok{"drop"}\NormalTok{)}
\FunctionTok{glimpse}\NormalTok{(panel)}
\end{Highlighting}
\end{Shaded}

\begin{verbatim}
## Rows: 168
## Columns: 6
## $ province      <chr> "Drenthe", "Drenthe", "Drenthe", "Drenthe", "Drenthe", "~
## $ year          <int> 2011, 2012, 2013, 2014, 2015, 2016, 2017, 2018, 2019, 20~
## $ avg_price     <dbl> 203389, 193844, 181850, 185239, 194566, 205504, 213562, ~
## $ median_income <dbl> 32000, 32100, 32200, 32900, 33200, 34200, 34800, 35600, ~
## $ pir           <dbl> 6.355906, 6.038754, 5.647516, 5.630365, 5.860422, 6.0088~
## $ gap           <dbl> 1.1186354, 1.0129180, 0.8825052, 0.8767882, 0.9534739, 1~
\end{verbatim}

\section{Visualisations}\label{visualisations}

\subsection{Temporal variation}\label{temporal-variation}

\begin{Shaded}
\begin{Highlighting}[]
\FunctionTok{ggplot}\NormalTok{(national, }\FunctionTok{aes}\NormalTok{(year, pir)) }\SpecialCharTok{+}
  \FunctionTok{geom\_hline}\NormalTok{(}\AttributeTok{yintercept =} \DecValTok{3}\NormalTok{, }\AttributeTok{linetype =} \StringTok{"dashed"}\NormalTok{, }\AttributeTok{colour =}\NormalTok{ vu\_navy) }\SpecialCharTok{+}
  \FunctionTok{geom\_line}\NormalTok{(}\AttributeTok{colour =}\NormalTok{ vu\_blue, }\AttributeTok{linewidth =} \DecValTok{1}\NormalTok{) }\SpecialCharTok{+} \FunctionTok{geom\_point}\NormalTok{(}\AttributeTok{colour =}\NormalTok{ vu\_navy) }\SpecialCharTok{+}
  \FunctionTok{annotate}\NormalTok{(}\StringTok{"text"}\NormalTok{, }\AttributeTok{x =} \FunctionTok{min}\NormalTok{(national}\SpecialCharTok{$}\NormalTok{year), }\AttributeTok{y =} \DecValTok{3}\NormalTok{,}
           \AttributeTok{label =} \StringTok{"UN{-}Habitat threshold = 3"}\NormalTok{, }\AttributeTok{vjust =} \SpecialCharTok{{-}}\FloatTok{0.6}\NormalTok{, }\AttributeTok{hjust =} \DecValTok{0}\NormalTok{,}
           \AttributeTok{colour =}\NormalTok{ vu\_navy, }\AttributeTok{size =} \DecValTok{3}\NormalTok{) }\SpecialCharTok{+}
  \FunctionTok{labs}\NormalTok{(}\AttributeTok{title =} \StringTok{"Affordability has steadily worsened"}\NormalTok{,}
       \AttributeTok{subtitle =} \FunctionTok{paste0}\NormalTok{(}\StringTok{"National average PIR, "}\NormalTok{, }\FunctionTok{min}\NormalTok{(national}\SpecialCharTok{$}\NormalTok{year), }\StringTok{"–"}\NormalTok{,}
                         \FunctionTok{max}\NormalTok{(national}\SpecialCharTok{$}\NormalTok{year)),}
       \AttributeTok{x =} \ConstantTok{NULL}\NormalTok{, }\AttributeTok{y =} \StringTok{"Price{-}to{-}income ratio"}\NormalTok{)}
\end{Highlighting}
\end{Shaded}

\begin{figure}

{\centering \includegraphics{report_files/figure-latex/fig-temporal-1} 

}

\caption{National price-to-income ratio over time.}\label{fig:fig-temporal}
\end{figure}

The PIR (years of median disposable income per average home) sits far
above the UN-Habitat affordability threshold of 3 throughout and trends
clearly upward, confirming a structural worsening of affordability over
the observed period.

\subsection{Sub-population variation}\label{sub-population-variation}

\begin{Shaded}
\begin{Highlighting}[]
\NormalTok{grp }\OtherTok{\textless{}{-}}\NormalTok{ panel }\SpecialCharTok{|\textgreater{}}
  \FunctionTok{group\_by}\NormalTok{(province) }\SpecialCharTok{|\textgreater{}} \FunctionTok{mutate}\NormalTok{(}\AttributeTok{mean\_inc =} \FunctionTok{mean}\NormalTok{(median\_income)) }\SpecialCharTok{|\textgreater{}} \FunctionTok{ungroup}\NormalTok{() }\SpecialCharTok{|\textgreater{}}
  \FunctionTok{mutate}\NormalTok{(}\AttributeTok{group =} \FunctionTok{ifelse}\NormalTok{(mean\_inc }\SpecialCharTok{\textgreater{}=} \FunctionTok{median}\NormalTok{(mean\_inc),}
                        \StringTok{"Higher{-}income provinces"}\NormalTok{, }\StringTok{"Lower{-}income provinces"}\NormalTok{))}
\FunctionTok{ggplot}\NormalTok{(grp, }\FunctionTok{aes}\NormalTok{(year, pir, }\AttributeTok{group =}\NormalTok{ province, }\AttributeTok{colour =}\NormalTok{ group)) }\SpecialCharTok{+}
  \FunctionTok{geom\_line}\NormalTok{(}\AttributeTok{alpha =}\NormalTok{ .}\DecValTok{5}\NormalTok{) }\SpecialCharTok{+}
  \FunctionTok{stat\_summary}\NormalTok{(}\FunctionTok{aes}\NormalTok{(}\AttributeTok{group =}\NormalTok{ group), }\AttributeTok{fun =}\NormalTok{ mean, }\AttributeTok{geom =} \StringTok{"line"}\NormalTok{, }\AttributeTok{linewidth =} \FloatTok{1.2}\NormalTok{) }\SpecialCharTok{+}
  \FunctionTok{facet\_wrap}\NormalTok{(}\SpecialCharTok{\textasciitilde{}}\NormalTok{group) }\SpecialCharTok{+}
  \FunctionTok{scale\_colour\_manual}\NormalTok{(}\AttributeTok{values =} \FunctionTok{c}\NormalTok{(vu\_blue, vu\_orange), }\AttributeTok{guide =} \StringTok{"none"}\NormalTok{) }\SpecialCharTok{+}
  \FunctionTok{labs}\NormalTok{(}\AttributeTok{title =} \StringTok{"Unaffordability is concentrated in richer provinces"}\NormalTok{,}
       \AttributeTok{subtitle =} \StringTok{"Thin lines = provinces; thick line = group mean"}\NormalTok{,}
       \AttributeTok{x =} \ConstantTok{NULL}\NormalTok{, }\AttributeTok{y =} \StringTok{"PIR"}\NormalTok{)}
\end{Highlighting}
\end{Shaded}

\begin{figure}

{\centering \includegraphics{report_files/figure-latex/fig-subpop-1} 

}

\caption{PIR for higher- vs. lower-income provinces.}\label{fig:fig-subpop}
\end{figure}

Higher-income (largely Randstad) provinces show both higher PIR levels
and faster growth, so the affordability problem is not uniform --- it is
regionally concentrated.

\subsection{Event association: the 2008 financial
crisis}\label{event-association-the-2008-financial-crisis}

\begin{Shaded}
\begin{Highlighting}[]
\NormalTok{event\_df }\OtherTok{\textless{}{-}}\NormalTok{ price\_history }\SpecialCharTok{|\textgreater{}}
  \FunctionTok{mutate}\NormalTok{(}\AttributeTok{period =} \FunctionTok{ifelse}\NormalTok{(year }\SpecialCharTok{\textless{}} \DecValTok{2009}\NormalTok{, }\StringTok{"Pre{-}2008"}\NormalTok{, }\StringTok{"2008 onward"}\NormalTok{))}
\FunctionTok{ggplot}\NormalTok{(event\_df, }\FunctionTok{aes}\NormalTok{(year, avg\_price, }\AttributeTok{colour =}\NormalTok{ period)) }\SpecialCharTok{+}
  \FunctionTok{geom\_vline}\NormalTok{(}\AttributeTok{xintercept =} \DecValTok{2008}\NormalTok{, }\AttributeTok{linetype =} \StringTok{"dotted"}\NormalTok{, }\AttributeTok{colour =}\NormalTok{ vu\_red) }\SpecialCharTok{+}
  \FunctionTok{geom\_line}\NormalTok{(}\FunctionTok{aes}\NormalTok{(}\AttributeTok{group =} \DecValTok{1}\NormalTok{), }\AttributeTok{linewidth =} \DecValTok{1}\NormalTok{) }\SpecialCharTok{+} \FunctionTok{geom\_point}\NormalTok{(}\AttributeTok{size =} \DecValTok{1}\NormalTok{) }\SpecialCharTok{+}
  \FunctionTok{annotate}\NormalTok{(}\StringTok{"text"}\NormalTok{, }\AttributeTok{x =} \DecValTok{2008}\NormalTok{, }\AttributeTok{y =} \FunctionTok{max}\NormalTok{(event\_df}\SpecialCharTok{$}\NormalTok{avg\_price),}
           \AttributeTok{label =} \StringTok{"2008 crisis"}\NormalTok{, }\AttributeTok{colour =}\NormalTok{ vu\_red, }\AttributeTok{hjust =} \SpecialCharTok{{-}}\FloatTok{0.05}\NormalTok{, }\AttributeTok{size =} \DecValTok{3}\NormalTok{) }\SpecialCharTok{+}
  \FunctionTok{scale\_y\_continuous}\NormalTok{(}\AttributeTok{labels =} \FunctionTok{label\_number}\NormalTok{(}\AttributeTok{prefix =} \StringTok{"€"}\NormalTok{, }\AttributeTok{big.mark =} \StringTok{","}\NormalTok{)) }\SpecialCharTok{+}
  \FunctionTok{scale\_colour\_manual}\NormalTok{(}\AttributeTok{values =} \FunctionTok{c}\NormalTok{(}\StringTok{"Pre{-}2008"} \OtherTok{=}\NormalTok{ vu\_grey, }\StringTok{"2008 onward"} \OtherTok{=}\NormalTok{ vu\_blue)) }\SpecialCharTok{+}
  \FunctionTok{labs}\NormalTok{(}\AttributeTok{title =} \StringTok{"The crisis interrupted, but did not reverse, the trend"}\NormalTok{,}
       \AttributeTok{x =} \ConstantTok{NULL}\NormalTok{, }\AttributeTok{y =} \StringTok{"Average price"}\NormalTok{, }\AttributeTok{colour =} \ConstantTok{NULL}\NormalTok{)}
\end{Highlighting}
\end{Shaded}

\begin{figure}

{\centering \includegraphics{report_files/figure-latex/fig-event-1} 

}

\caption{Average house price around the 2008 crisis (real price data, 1995-2024).}\label{fig:fig-event}
\end{figure}

Prices fell for roughly five years after 2008 before resuming their
climb, showing the crisis as a temporary correction rather than a
structural break.

\subsection{Spatial variation: a choropleth
map}\label{spatial-variation-a-choropleth-map}

\begin{Shaded}
\begin{Highlighting}[]
\ControlFlowTok{if}\NormalTok{ (}\SpecialCharTok{!}\FunctionTok{is.null}\NormalTok{(nl) }\SpecialCharTok{\&\&} \StringTok{"statcode"} \SpecialCharTok{\%in\%} \FunctionTok{names}\NormalTok{(nl)) \{}
\NormalTok{  nl2 }\OtherTok{\textless{}{-}}\NormalTok{ nl }\SpecialCharTok{|\textgreater{}} \FunctionTok{left\_join}\NormalTok{(pir\_latest, }\AttributeTok{by =} \StringTok{"statcode"}\NormalTok{)}
  \FunctionTok{ggplot}\NormalTok{(nl2) }\SpecialCharTok{+}
    \FunctionTok{geom\_sf}\NormalTok{(}\FunctionTok{aes}\NormalTok{(}\AttributeTok{fill =}\NormalTok{ pir), }\AttributeTok{colour =} \StringTok{"white"}\NormalTok{, }\AttributeTok{linewidth =}\NormalTok{ .}\DecValTok{2}\NormalTok{) }\SpecialCharTok{+}
    \FunctionTok{scale\_fill\_gradient}\NormalTok{(}\AttributeTok{low =} \StringTok{"\#cfe8f5"}\NormalTok{, }\AttributeTok{high =}\NormalTok{ vu\_red, }\AttributeTok{name =} \StringTok{"PIR"}\NormalTok{) }\SpecialCharTok{+}
    \FunctionTok{labs}\NormalTok{(}\AttributeTok{title =} \FunctionTok{paste0}\NormalTok{(}\StringTok{"Least affordable provinces, "}\NormalTok{, latest),}
         \AttributeTok{subtitle =} \StringTok{"Darker = higher price{-}to{-}income ratio"}\NormalTok{,}
         \AttributeTok{caption =} \StringTok{"Source: CBS 83625NED, 86161NED; boundaries cartomap.github.io"}\NormalTok{) }\SpecialCharTok{+}
    \FunctionTok{theme\_void}\NormalTok{(}\AttributeTok{base\_size =} \DecValTok{11}\NormalTok{) }\SpecialCharTok{+} \FunctionTok{theme}\NormalTok{(}\AttributeTok{legend.position =} \StringTok{"right"}\NormalTok{)}
\NormalTok{\} }\ControlFlowTok{else}\NormalTok{ \{}
  \CommentTok{\# Fallback: ranked bar chart if the boundary file can\textquotesingle{}t be fetched offline.}
  \FunctionTok{ggplot}\NormalTok{(pir\_latest, }\FunctionTok{aes}\NormalTok{(}\FunctionTok{reorder}\NormalTok{(province, pir), pir, }\AttributeTok{fill =}\NormalTok{ pir)) }\SpecialCharTok{+}
    \FunctionTok{geom\_col}\NormalTok{() }\SpecialCharTok{+} \FunctionTok{coord\_flip}\NormalTok{() }\SpecialCharTok{+}
    \FunctionTok{scale\_fill\_gradient}\NormalTok{(}\AttributeTok{low =} \StringTok{"\#cfe8f5"}\NormalTok{, }\AttributeTok{high =}\NormalTok{ vu\_red, }\AttributeTok{guide =} \StringTok{"none"}\NormalTok{) }\SpecialCharTok{+}
    \FunctionTok{labs}\NormalTok{(}\AttributeTok{title =} \FunctionTok{paste0}\NormalTok{(}\StringTok{"PIR by province, "}\NormalTok{, latest),}
         \AttributeTok{subtitle =} \StringTok{"Map boundaries unavailable offline — ranked bars shown"}\NormalTok{,}
         \AttributeTok{x =} \ConstantTok{NULL}\NormalTok{, }\AttributeTok{y =} \StringTok{"PIR"}\NormalTok{)}
\NormalTok{\}}
\end{Highlighting}
\end{Shaded}

\begin{figure}

{\centering \includegraphics{report_files/figure-latex/fig-map-1} 

}

\caption{Price-to-income ratio by province (latest year).}\label{fig:fig-map}
\end{figure}

The map shows a clear Randstad-versus-periphery divide: the western
provinces carry the highest price-to-income ratios, confirming that the
affordability problem is spatially concentrated rather than national in
scale.

\section{Key findings}\label{key-findings}

Nationally the PIR rose from 7.1** in 2011 to 8.8** in 2024. In 2024 the
least affordable province was Noord-Holland** (PIR 11.5) and the most
affordable Zeeland** (PIR 7.3) --- a gap of 4.3** points, underscoring
the regional dimension.

\section{Limitations}\label{limitations}

National and provincial averages mask within-province variation (cities
vs. rural). The analysis covers owner-occupied purchases only, excluding
the large rental and social-housing sector that binds for many
households. Figures are nominal unless stated, so part of the price
growth reflects inflation. Because the regional income series starts in
2011, the affordability panel spans 2011-2024; the 2008 event analysis
therefore uses the price series alone. The event analysis is descriptive
and cannot establish causality. Finally, the PIR uses disposable
(post-tax) income, so its levels sit above the gross-income convention
behind the UN-Habitat ``3'' benchmark and are best read as a relative
indicator across provinces and time.

\section{References}\label{references}

\textbf{Literature (problem motivation)}

\begin{itemize}
\tightlist
\item
  OECD (2025). \emph{OECD Economic Surveys: Netherlands 2025 --- Towards
  a more accessible and sustainable housing market.} OECD Publishing.
  \url{https://www.oecd.org/en/publications/oecd-economic-surveys-netherlands-2025_2dd1f4aa-en.html}
\item
  De Nederlandsche Bank (2024). \emph{Borrowing capacity and house
  prices} (DNB analysis).
  \url{https://www.dnb.nl/media/ccsjet4s/dnb-analysis-borrowing-capacity-and-house.pdf}
\item
  UN-Habitat (n.d.). \emph{Housing affordability and the price-to-income
  ratio.} \url{https://unhabitat.org}
\end{itemize}

\textbf{Data sources}

\begin{itemize}
\tightlist
\item
  Statistics Netherlands (CBS). \emph{Bestaande koopwoningen; prijzen,
  regio} (83625NED) and \emph{Inkomen van huishoudens; regio}
  (86161NED), via the \texttt{cbsodataR} package.
  \url{https://opendata.cbs.nl}
\item
  Province boundaries: cartomap (\url{https://cartomap.github.io/nl}).
\end{itemize}

\textbf{Software}

\begin{itemize}
\tightlist
\item
  Wickham, H. et al.~(2019). \emph{Welcome to the tidyverse.} JOSS
  4(43), 1686.
\end{itemize}

\end{document}
